A judge that follows position alone picks a different image in each rotation, so it can never be unanimous. Unanimity keeps 121 prompts (40.3\%). On them Qwen's gain over random is 0.141 [0.110, 0.173] and regret 0.074; on the 179 rejected prompts its gain is $-$0.030 [$-$0.057, $-$0.002], below random (Figure~\ref{fig:results}b). The kept prompts are not easier: their candidates differ more in rating than the rejected ones (mean best-minus-worst range 0.454 versus 0.404), leaving more room for error. The gate only filters prompts; on kept prompts the choice is unchanged.

The two-of-three gate is an intermediate operating point: it keeps 80.7\% of prompts with gain 0.094 [0.071, 0.118] (Appendix Table~\ref{tab:policies}).

Voting over the three orders instead of abstaining, a simple form of order averaging, raises the gain on all prompts from 0.039 to 0.065 [0.043, 0.086].

Agreement between judges behaves differently. Qwen and Smol agree on 59 prompts (19.7\%), where Qwen's gain is $-$0.060 [$-$0.111, $-$0.011]; on the 241 prompts where they disagree it is 0.064 [0.039, 0.087]. Two judges that both prefer early slots can agree for reasons unrelated to content, so this gate keeps the wrong prompts. The result is specific to pairing with a weak judge. Appendix Table~\ref{tab:policies} lists every gate with its complement.
